\documentclass[10pt,a4paper]{amsart}
\usepackage{amsmath,amssymb,amsthm,graphicx}
\usepackage{verbatim}
\newenvironment{reference}{\comment}{\endcomment}
%\newenvironment{reference}{}{}
\newenvironment{slogan}{\comment}{\endcomment}
\newenvironment{history}{\comment}{\endcomment}

% For commutative diagrams we use Xy-pic
\usepackage[all]{xy}

% We use 2cell for 2-commutative diagrams.
\xyoption{2cell}
% Theorem environments.
%
\theoremstyle{plain}
\newtheorem{theorem}[subsection]{Theorem}
\newtheorem{proposition}[subsection]{Proposition}
\newtheorem{lemma}[subsection]{Lemma}

\theoremstyle{definition}
\newtheorem{definition}[subsection]{Definition}
\newtheorem{example}[subsection]{Example}
\newtheorem{exercise}[subsection]{Exercise}
\newtheorem{situation}[subsection]{Situation}

\theoremstyle{remark}
\newtheorem{remark}[subsection]{Remark}
\newtheorem{remarks}[subsection]{Remarks}

\numberwithin{equation}{subsection}

% Macros
%
\def\lim{\mathop{\mathrm{lim}}\nolimits}
\def\colim{\mathop{\mathrm{colim}}\nolimits}
\def\Spec{\mathop{\mathrm{Spec}}}
\def\Hom{\mathop{\mathrm{Hom}}\nolimits}
\def\Ext{\mathop{\mathrm{Ext}}\nolimits}
\def\SheafHom{\mathop{\mathcal{H}\!\mathit{om}}\nolimits}
\def\SheafExt{\mathop{\mathcal{E}\!\mathit{xt}}\nolimits}
\def\Sch{\mathit{Sch}}
\def\Mor{\mathop{\mathrm{Mor}}\nolimits}
\def\Ob{\mathop{\mathrm{Ob}}\nolimits}
\def\Sh{\mathop{\mathit{Sh}}\nolimits}
\def\NL{\mathop{N\!L}\nolimits}
\def\CH{\mathop{\mathrm{CH}}\nolimits}
\def\proetale{{pro\text{-}\acute{e}tale}}
\def\etale{{\acute{e}tale}}
\def\QCoh{\mathit{QCoh}}
\def\Ker{\mathop{\mathrm{Ker}}}
\def\Im{\mathop{\mathrm{Im}}}
\def\Coker{\mathop{\mathrm{Coker}}}
\def\Coim{\mathop{\mathrm{Coim}}}

% Boxtimes
%
\DeclareMathSymbol{\boxtimes}{\mathbin}{AMSa}{"02}

%
% Macros for moduli stacks/spaces
%
\def\QCohstack{\mathcal{QC}\!\mathit{oh}}
\def\Cohstack{\mathcal{C}\!\mathit{oh}}
\def\Spacesstack{\mathcal{S}\!\mathit{paces}}
\def\Quotfunctor{\mathrm{Quot}}
\def\Hilbfunctor{\mathrm{Hilb}}
\def\Curvesstack{\mathcal{C}\!\mathit{urves}}
\def\Polarizedstack{\mathcal{P}\!\mathit{olarized}}
\def\Complexesstack{\mathcal{C}\!\mathit{omplexes}}
% \Pic is the operator that assigns to X its picard group, usage \Pic(X)
% \Picardstack_{X/B} denotes the Picard stack of X over B
% \Picardfunctor_{X/B} denotes the Picard functor of X over B
\def\Pic{\mathop{\mathrm{Pic}}\nolimits}
\def\Picardstack{\mathcal{P}\!\mathit{ic}}
\def\Picardfunctor{\mathrm{Pic}}
\def\Deformationcategory{\mathcal{D}\!\mathit{ef}}
\usepackage{cleveref}
\crefname{originallemma}{original lemma}{original lemmas}
\Crefname{originallemma}{Original lemma}{Original lemmas}
\makeatletter
\newlabel{lemma-Noether-normalization}{{00OY}{source}{}{sourcepointer-1}{}}
\newlabel{lemma-Noether-normalization@cref}{{[originallemma][0][]00OY}{[1][0][]source}{}{}{}}
\makeatother
\crefname{originallemma}{original lemma}{original lemmas}
\Crefname{originallemma}{Original lemma}{Original lemmas}
\makeatletter
\newlabel{lemma-characterize-finite-in-terms-of-integral}{{02JJ}{source}{}{sourcepointer-2}{}}
\newlabel{lemma-characterize-finite-in-terms-of-integral@cref}{{[originallemma][0][]02JJ}{[1][0][]source}{}{}{}}
\makeatother
\crefname{originallemma}{original lemma}{original lemmas}
\Crefname{originallemma}{Original lemma}{Original lemmas}
\makeatletter
\newlabel{lemma-finite-is-integral}{{00GK}{source}{}{sourcepointer-3}{}}
\newlabel{lemma-finite-is-integral@cref}{{[originallemma][0][]00GK}{[1][0][]source}{}{}{}}
\makeatother
\title{Noether normalization over a domain, up to inverting one element of the base}
\author{The Stacks Project Authors}
\date{stacks-project:algebra.tex:local-snapshot-20261012; GNU FDL 1.2 or later}
\begin{document}
\maketitle
\noindent\textit{Editorial note.} All delivered mathematical text is the author's own, copied from the pinned Stacks Project chapter file algebra.tex; no statement or proof step was rewritten, reordered or removed.\par\medskip
\noindent\textit{Editorial note.} Preamble provenance, as required for this source: the frozen snapshot stages the chapter's own preamble file beside it (preamble.tex, 217 lines, 8116 bytes, the same shared Stacks-wide preamble the published chow and etale-cohomology cards use); the first line of the pinned chapter file is \textbackslash{}input\{preamble\}, so no preamble text lives inside the chapter itself. The card takes that file's lines 8-18 (the verbatim package and the three \textbackslash{}comment-based environments reference/slogan/history at lines 9-12, section 13.2, together with the Xy-pic setup) and its lines 154 to the runtime end of the file (the native theorem declarations and mathematical macros). Not one character was changed, no package was added, no macro was defined, and the excerpt serves only as preamble, never as a span coordinate.\par\medskip
\noindent\textit{Editorial note.} The same-chapter pointers are located at their \textbackslash{}label lines inside the pinned chapter file; the cross-chapter pointers are located in the project's own tag list staged in the snapshot (sources/stacks-chow/tags.txt), at the verbatim lines carrying their canonical tags.\par\medskip
\noindent\textit{Editorial note.} The counted proof refers to the absolute normalization statement by label; it is recorded as a navigation stub pointed at its \textbackslash{}label line, and the argument stays inside the counted span.\par\medskip
\noindent\textit{Editorial note.} The counted spans contain no \textbackslash{}cite, so the card carries no bibliography.\par\medskip
\noindent\textit{Allowed standard prerequisites (editorial).}\begin{itemize}
\item Noether normalization over the fraction field. Invoked at the start of the counted proof (lines 28355-28360). Produces the algebraically independent elements over which the generic fibre is finite.
\item Finiteness spread out by clearing finitely many denominators. Used in the second half of the counted proof (lines 28360-28377). Turns the finiteness over the fraction field into finiteness after inverting one element of the base.
\end{itemize}\medskip
\ifdefined\hypertarget\hypertarget{sourcepointer-1}{}\fi
\noindent\textit{Source reference.} Original lemma 00OY: sources/stacks-algebra/pinned-original-source.tex, source lines 28242--28242. Navigation pointer for \textbackslash{}ref\{lemma-Noether-normalization\} used inside the counted span (line 28356); the argument that consumes it is retained in the span.\par\medskip
\ifdefined\hypertarget\hypertarget{sourcepointer-2}{}\fi
\noindent\textit{Source reference.} Original lemma 02JJ: sources/stacks-algebra/pinned-original-source.tex, source lines 7720--7720. Navigation pointer for \textbackslash{}ref\{lemma-characterize-finite-in-terms-of-integral\} used inside the counted span (line 28373); the argument that consumes it is retained in the span.\par\medskip
\ifdefined\hypertarget\hypertarget{sourcepointer-3}{}\fi
\noindent\textit{Source reference.} Original lemma 00GK: sources/stacks-algebra/pinned-original-source.tex, source lines 7687--7687. Navigation pointer for \textbackslash{}ref\{lemma-finite-is-integral\} used inside the counted span (line 28361); the argument that consumes it is retained in the span.\par\medskip
\begin{lemma}
\label{lemma-Noether-normalization-over-a-domain}
Let $R \to S$ be an injective finite type ring map. Assume $R$ is a domain.
Then there exists an integer $d$ and a factorization
$$
R \to R[y_1, \ldots, y_d] \to S' \to S
$$
by injective maps such that $S'$ is finite over $R[y_1, \ldots, y_d]$
and such that $S'_f \cong S_f$ for some nonzero $f \in R$.
\end{lemma}

\begin{proof}
Pick $x_1, \ldots, x_n \in S$ which generate $S$ over $R$.
Let $K$ be the fraction field of $R$ and $S_K = S \otimes_R K$.
By Lemma \ref{lemma-Noether-normalization}
we can find $y_1, \ldots, y_d \in S$ such that $K[y_1, \ldots, y_d] \to S_K$
is a finite injective map. Note that $y_i \in S$ because we may pick the
$y_j$ in the $\mathbf{Z}$-algebra generated by $x_1, \ldots, x_n$.
As a finite ring map is integral (see
Lemma \ref{lemma-finite-is-integral})
we can find monic $P_i \in K[y_1, \ldots, y_d][T]$ such that
$P_i(x_i) = 0$ in $S_K$. Let $f \in R$ be a nonzero element such that
$fP_i \in R[y_1, \ldots, y_d][T]$ for all $i$. Then $fP_i(x_i)$
maps to zero in $S_K$. Hence after replacing $f$ by another
nonzero element of $R$ we may also assume $fP_i(x_i)$ is zero in $S$.
Set $x_i' = fx_i$ and let $S' \subset S$ be the $R$-subalgebra generated by
$y_1, \ldots, y_d$ and $x'_1, \ldots, x'_n$. Note that $x'_i$ is integral over
$R[y_1, \ldots, y_d]$ as we have $Q_i(x_i') = 0$ where
$Q_i = f^{\deg_T(P_i)}P_i(T/f)$ which is a monic polynomial in
$T$ with coefficients in $R[y_1, \ldots, y_d]$ by our choice of $f$.
Hence $R[y_1, \ldots, y_d] \subset S'$ is finite by
Lemma \ref{lemma-characterize-finite-in-terms-of-integral}.
Since $S' \subset S$ we have $S'_f \subset S_f$ (localization is exact).
On the other hand, the elements $x_i = x'_i/f$ in $S'_f$ generate $S_f$
over $R_f$ and hence $S'_f \to S_f$ is surjective. Whence
$S'_f \cong S_f$ and we win.
\end{proof}

\end{document}
